% =============================================================================
% REFERANDUM, Nisan 2017
% =============================================================================
\newpage
\section{Kimlikler}

Hem ideolojik hem de etnik ve inanç kimlikleri olaylara duyarlı. 
Etnik ve inanç kimlikleri Özgür Özel döneminde söylemde yer ediniyor. CHP söylemi ideolojik kimliklerden çok etnik
ve inanç kimliklerini yöneliyor. Etnik ve inanç kimlikleri bütün olaylara tepki gösterirken ideolojik kimlikler referandum, 2019 yerel seçimleri ve İmamoğlu'nun
tutuklanması gibi dislokasyonlara duyarsız kalıyor. 


\begin{table}[H]
\centering
\caption{Trendlerine göre ideolojik ile etnik ve inanç kimlikleri kayıtları}
\label{tab:olay-kayit}
\small
\begin{tabular}{@{}l l r r r c l@{}}
\toprule
Kayıt & Olay & Taban & Sonrası & Değişim & & Dönüş \\
\midrule

% ---------------------------------------------------------------- ÇEKİLENLER
\multicolumn{7}{@{}l}{\textbf{Çekilenler}} \\[0.2em]
Etnik ve inanç & Darbe girişimi, 2016 & 2.3 & 0.6 & $-$72\% & \azalan & > 24 ay \\
Etnik ve inanç & Referandum, 2017 & 0.6 & 0.4 & $-$42\% & \azalan & 4 ay \\
Etnik ve inanç & Genel seçim, 2018 & 1.0 & 0.0 & $-$95\% & \azalan & 15 ay \\
İdeolojik & Genel seçim, 2018 & 4.9 & 3.3 & $-$33\% & \azalan & 3 ay \\
Etnik ve inanç & Genel seçim, 2023 & 1.1 & 0.2 & $-$85\% & \azalan & 5 ay \\
İdeolojik & Bahçeli çağrısı, 2024 & 7.0 & 4.7 & $-$32\% & \azalan & 4 ay \\
İdeolojik & İstanbul kayyumu, 2025 & 5.2 & 2.8 & $-$46\% & \azalan & gözlem dışı \\

\addlinespace[0.6em]
% ---------------------------------------------------------------- YÜKSELENLER
\multicolumn{7}{@{}l}{\textbf{Yükselenler}} \\[0.2em]
İdeolojik & Darbe girişimi, 2016 & 0.6 & 2.3 & $+$283\% & \artan \\
Etnik ve inanç & Yerel seçim, 2019 & 0.3 & 0.8 & $+$142\% & \artan \\
Etnik ve inanç & Yerel seçim, 2024 & 2.3 & 5.6 & $+$138\% & \artan \\
Etnik ve inanç & Bahçeli çağrısı, 2024 & 5.1 & 6.4 & $+25$\% & \artan \\
Etnik ve inanç & İmamoğlu tutuklanması, 2025 & 6.0 & 9.8 & $+$62\% & \artan \\
Etnik ve inanç & İstanbul kayyumu, 2025 & 8.0 & 8.9 & $+$12\% & \artan \\


\addlinespace[0.6em]
% ------------------------------------------------------------ HAREKET ETMEYEN
\multicolumn{7}{@{}l}{\textbf{Hareket etmeyenler}} \\[0.2em]
İdeolojik & Referandum, 2017 & 3.6 & 4.0 & $+$10\% &  &  \\
İdeolojik & Yerel seçim, 2019 & 5.1 & 5.1 & 0\% & & \\
İdeolojik & İmamoğlu tutuklaması, 2025 & 5.8 & 6.0 & $+$3\% & & \\

\bottomrule
\end{tabular}
\kaynaknot{Taban olaydan önceki 12 ayın, \emph{Sonrası} olaydan sonraki
üç ayın ortalamasıdır; ikisi de 10.000 kelimede vukuat sayısıdır, yani
konuşma hacminden bağımsızdır. \emph{Dönüş}, kaydın tabanın \%90'ına geri
döndüğü ay sayısıdır. Korpus 31 Mart 2026'da bittiği için Bahçeli çağrısı,
İmamoğlu tutuklaması ve İstanbul kayyumu sonrasında 24 aylık pencere
doldurulamıyor; oralarda dönüş \emph{gözlem dışıdır}, gerçekleşmediği
anlamına gelmez.}
\end{table}

\subsection{Kadın: Semantik}
\begin{figure}[H]
    \centering
    \includegraphics[width=\textwidth]{kadın_phase_heatmap.pdf}
    \caption[Phase heatmap: \emph{kadın}]{Kadın kelimesi etrafındaki bağlam kümelerinin
    faz oranı. Oran her fazdaki bir kümenin payının tüm korpustaki payına oranıdır. 
    1.0 değeri orantılı varlığı; 1.0 üzeri veya altı değerler göreli yoğunluk veya kıtlığı gösterir.
    Renk koyulaştıkça ilgili kümenin o fazdaki yoğunluğu artar. $KK_{Erken} = 13.01.2015-01.04.2019$; 
    $KK_{Gec} = 02.04.2019-09.11.2023$; $Ozel = 06.11.2023-31.03.2026$. Sırasıyla dönem frekansları: 1: 1161, 2: 1884, 3: 2563.}
    \label{fig:heatmap-kadın}
\end{figure}

Kılıçdaroğlu, erken dönemde kadınlara yönelik söylemini popüler taleplerle birlikte dile getirdiği kadın-erkek eşitliği
ve genç, yaşlı, erkek ve çocuğun beraber-kapsayıcı mücadelesi üzerinden inşa etmektedir. Okuryazarlık ve Eğitim ile Seçme
ve Seçilme kümeleri bir arada değerlendirildiğinde Kemal Kılıçdaroğlu erken Cumhuriyet döneminin 
kazanımlarına daha yoğun atıfta bulunmaktadır. Bahşedilen Haklar kümesinde Kılıçdaroğlu,
"Cennet annelerin ayakları altındadır." hadisi ve çeşitli dini referanslar\footnote{Benzeri retorik dini referanslara tezimde "hak" kelimesinin etrafındaki semantik kümeleri incelerken de rastladım. Örn: Kul hakkı, yetimin ve öksüzün hakkı gibi tamlamalar.}
üzerinden kadına şiddeti ve AKP'nın önlemsizliğini eleştirmekte, aynı kümede
Özgür Özel ise dini referansları bırakıp yerine, Kılıçdaroğlu'nda da gözlemlendiği gibi, erken Cumhuriyet döneminde kadınlara tanınan haklardan söz etmektedir.

Geç dönem söz konusu olduğunda Kılıçdaroğlu'nun söylemindeki Atatürk döneminin getirilerine yönelik anlatı
sönümlenmiş ve önceki dönemden aktif olan 8 kümeden yalnızca Cinsiyet Eşitliği, kapsayıcı söylem, Kadınların Siyasete Katılımı ve ahlaki 
retoriği devam etmiştir. Bu dönemde Kılıçdaroğlu'nun sosyal yardımlara yönelik vaatleri ilgi çekici olup (örn. Aile Destekleri Sigortası ve Sosyal Güvenlik Primi\footnote{Bu vaatten ilgili kümede ilk kez 3 Eylül 2021 tarihinde bahsediyor.}) ağırlıklı olarak 2023 Mayıs ayındaki mitinglerde yoğunlaşmaktadır.
2020 yılının Temmuz ayından başlamak itibariyle CHP lidelerinin söylemine kadın örgütleriyle yapılan
görüşmeler de konu olmaya başlamıştır. Kılıçdaroğlu, 306 kadın örgütüyle yazılmış yeni siyasi partiler yasasını ve bu yasadaki
50-50 kadın kota önerisini sıklıkla dile getirirken Özgür Özel kadına yönelik şiddet, 6284 yasası ve Etki Ajanlığı Yasası hakkında
eleştiriler gündeme getirmektedir.

Isı haritasının orta kısmında yer alan kümeler CHP'nin son 11 yılda kadınlar hakkındaki
söyleminin temel bileşenleridir. Bu kümeler sırasıyla şu şekildedir: Kadın tutuklular ve mağdurlar, İstanbul Sözleşmesi ve kadına yönelik şiddet,
CHP belediyelerinin kreş girişimleri ve sosyal adaletsizlik, ve sosyal kimlikler. "Mesleki Eşdeğerlilik" anlatısı altında Kılıçdaroğlu işçi, çiftçi, emekli ve esnaf gibi 
"çalışan" sınıfları hem ekonomik problemlerle yüzleşen hem de ortak bir dayanışma içinde olan "sosyal kimlikler" olarak eklemlemektedir.
Ayrıca bütün bu kimliklerin çözümünün CHP'de olduğu, Bay Kemal'in bu kimlikler için çalışacağı gibi retorik izler de kümede mevcuttur.

Her ne kadar ağırlığı Kılıçdaroğlu'nun ikinci döneminde olsa da 
"Çalışma" kümesi de bu listeye dahil edilebilir. Kılıçdaroğlu, özellikle kırsaldaki kadınların ev içi emeğinin karşılıksız kalmasını eleştiriyorken
Özgür Özel, İBB tarafından açılan kreşlerin kadınların sosyal hayatına etkisine değinmektedir.

Özel döneminde partinin kadın söylemi AKP ve Erdoğan'a eklemlenmeye başlamıştır. Sırasıyla, kabinedeki tek bakanın Aile ve Sosyal Hizmetler
Bakanı olması, Erdoğan ailesine yönelik İran rejimindeki "kadın haklarına" yönelik ihlaller, Filistin'de kadın ve çocuklara yönelik uygulanan katliam,
Suriyeli kadın göçmenlerin yaşadığı zorluklar, MHP ve AKP'nin HÜDAPAR ile ittifak yapması kadınlara yönelik son yıllardaki CHP gündeminin
temel gündemlerini oluşturmaktadır. Buna ek olarak Özgür Özel döneminde gençler ve kadınlara yönelik ortak siyasa önerileri (örn: kota) 
söylemde daha sık yer edinmeye başlamıştır. İlginç bir biçimde Özgür Özel AKP'ye oy veren kadın seçmenlerin çocuklarının aldığı eğitimden memnuniyet oranı Özgür Özel'in
söyleminde sürekli gündeme gelmektedir. 2024 Eylül ayından başlayarak 2026 Ocak ayının başına kadar bu verinin \%17 ila \%21 arasında olduğu tarafınca yinelenmiştir\footnote{Bu veriyi nereden aldınız, güvenemedim.}. 

Özetle, 
\vspace{-1.2em}
\begin{itemize}
    \item Özgür Özel döneminde CHP'nin kadın söylemi daha politize olmuştur. Bunun temel sebebi Erdoğan'a, ailesine, AKP'ye ve bakanlığa yönelik direkt ve 
    kişisel eleştiriler, yani karşı bir düşman inşası, söyleme girmesidir\footnote{Yalnızca "kadın" etrafındaki semantik analizden çıkarılan sonuçtur, diğer kümeler ve frekanslar beraber değerlendirilmeden böyle bir sonuca varılamaz.}. Kemal Kılıçdaroğlu'nun kadınlara yönelik söylemi sosyal yardımlara öncelik veren ideal bir anlatıya bağlı kalırken
    Özgür Özel işsizlik, yoksuzlluk, kreş, kota uygulamaları gibi güncel ve siyasa üretmeye yönelik talepleri vurgulamaktadır.
    \item 
\end{itemize}

\subsection{Genç: Semantik}

% burada "zincir" kullandım, kadın kümelerindeki çalışma için de Zincir: Çalışma yazabilirim
 
